%==============================================================================
%  Supplementary File S3 — SQL evaluation metrics
%  Companion to main_applsci.tex. Compiles standalone; export to DOCX for submission.
%  The 72 expert-rating items referenced in Methods > SQL Query Validation.
%==============================================================================
\documentclass[10pt]{article}
\usepackage[utf8]{inputenc}
\usepackage[T1]{fontenc}
\usepackage[margin=0.9in]{geometry}
\usepackage{times}
\usepackage{booktabs}
\usepackage{longtable}
\usepackage{array}
\usepackage{amsmath}
\usepackage{amssymb}
\usepackage{xcolor}
\usepackage{caption}
\usepackage{enumitem}
\usepackage[colorlinks=true,allcolors=blue!55!black]{hyperref}

\newcommand{\PH}[1]{{\color{red}\textbf{[#1]}}}
\newcolumntype{L}[1]{>{\raggedright\arraybackslash}p{#1}}
\newcolumntype{C}[1]{>{\centering\arraybackslash}p{#1}}
\captionsetup{labelsep=period,font=small,justification=justified,
              singlelinecheck=false,skip=4pt}
\setcounter{secnumdepth}{0}
\pagestyle{plain}

\begin{document}

\begin{center}
{\large\bfseries Supplementary File S3. SQL evaluation metrics}\\[2pt]
{\footnotesize Companion to: \emph{Metadata-Grounded Large Language Models for Converting Hospital Statistical
Requests into Clinical Data Warehouse Queries}}
\end{center}

\vspace{6pt}
\noindent
Two experts independently rated every generated query against the 72 items
below. Each item is scored on a 4-point scale: 1 = noncompliant, 2 = partially
compliant with a defect that changes the result set, 3 = compliant with a defect
that does not change the result set, 4 = fully compliant. An item that does not
apply to a given query is scored NA and excluded from that query's mean.
Dimension scores are the unweighted mean of the applicable items; the reported
score is the mean of the two experts. Inter-rater reliability was computed with
Cohen $\kappa$ over the dichotomized ratings ($\ge$3 vs $<$3), and disagreements
were resolved by consensus discussion.

\vspace{4pt}
\noindent
The column \textbf{Subset} marks the 18 items that constitute the restricted
instrument applied to queries containing prevalidated criteria from the
development set; those 18 items measure \emph{criterion inclusion accuracy}
(12 items, prefix X) and \emph{field identifier correctness} (6 items,
prefix C), reported separately in the Results.

%------------------------------------------------------------------
\begingroup\footnotesize
\begin{longtable}{@{}L{1.0cm}L{12.4cm}C{1.4cm}@{}}
\caption{The 72 expert-rating items, by dimension.}\\
\toprule
ID & Item & Subset\\
\midrule
\endfirsthead
\multicolumn{3}{l}{\footnotesize\emph{Table continued from previous page}}\\
\toprule
ID & Item & Subset\\
\midrule
\endhead
\bottomrule
\multicolumn{3}{r}{\footnotesize\emph{Continued on next page}}\\
\endfoot
\bottomrule
\endlastfoot
\multicolumn{3}{@{}l}{\textbf{Dimension 1 --- SQL syntax adherence (24 items)}}\\
S1  & Statement parses under the Presto grammar & \\
S2  & Exactly one statement is returned, with no trailing fragment & \\
S3  & Every referenced alias is defined & \\
S4  & No duplicate alias within a scope & \\
S5  & Every SELECT expression is aggregated or appears in GROUP BY & \\
S6  & GROUP BY references valid expressions & \\
S7  & ORDER BY references valid output columns & \\
S8  & No reserved word used as an unquoted identifier & \\
S9  & Non-ASCII identifiers are enclosed in double quotes & \\
S10 & String literals use single quotes & \\
S11 & No implicit comparison between a VARCHAR column and a date or number & \\
S12 & Explicit cast or parse function before every date comparison & \\
S13 & Explicit cast before every numeric comparison & \\
S14 & Presto date function names used correctly & \\
S15 & No function specific to another SQL dialect & \\
S16 & Interval arithmetic uses Presto syntax & \\
S17 & NULL tested with IS NULL / IS NOT NULL, never with equality & \\
S18 & Division is guarded against a zero denominator & \\
S19 & COUNT(DISTINCT $\cdot$) used wherever deduplication is required & \\
S20 & Subqueries are correlated correctly & \\
S21 & Common table expressions are uniquely named and referenced & \\
S22 & No SELECT * in the returned statement & \\
S23 & Output aliases are present and non-colliding & \\
S24 & Statement executes without runtime error against the target instance & \\
\addlinespace

\multicolumn{3}{@{}l}{\textbf{Dimension 2 --- Warehouse schema compliance (26 items)}}\\
C1  & Every referenced table exists in the catalog & \checkmark\\
C2  & Every referenced column exists on its stated table & \checkmark\\
C3  & Each column is qualified with the table that actually owns it & \checkmark\\
C4  & Correct catalog and schema prefix for the subject area & \\
C5  & Layer choice is justified; a DC-layer aggregate is used when one exists & \\
C6  & No cross-layer join that double counts the same fact & \\
C7  & Join keys are the documented keys for the pair & \\
C8  & Join cardinality does not inflate the row set & \\
C9  & No join that is unnecessary for the requested output and can only reduce rows & \\
C10 & Every joined table carries its logical-deletion filter & \\
C11 & Every multi-campus table carries its organization-code filter & \\
C12 & Filters are placed in WHERE or ON correctly with respect to outer joins & \\
C13 & Correct main-table pairing for the subject area & \\
C14 & Required business-status filters are present & \\
C15 & Coded columns are compared against their declared value domain & \\
C16 & Code values are drawn from the value dictionary and not invented & \\
C17 & Units are consistent with the column definition & \\
C18 & Date columns are handled in their documented storage format & \\
C19 & No deprecated column is used & \\
C20 & No column with a population rate below threshold is used without a note & \\
C21 & Partition or pruning columns are used where available & \\
C22 & Field provenance is emitted for every selected column & \checkmark\\
C23 & Provenance names the table, the layer, and the mapping entry & \checkmark\\
C24 & No column is referenced that is absent from the resolved mapping & \checkmark\\
C25 & Text comparisons normalize casing and whitespace where required & \\
C26 & Output granularity matches the declared grouping & \\
\addlinespace

\multicolumn{3}{@{}l}{\textbf{Dimension 3 --- Criteria contextual accuracy (22 items)}}\\
X1  & The counting object matches the requested indicator & \checkmark\\
X2  & The time criterion matches the metric-criterion dictionary entry & \checkmark\\
X3  & No timestamp is silently substituted for the requested one & \checkmark\\
X4  & The reporting interval is the requested closed interval & \\
X5  & Relative dates are resolved against the stated report date & \\
X6  & Inclusion conditions are complete & \checkmark\\
X7  & Exclusion conditions are complete and correctly scoped & \checkmark\\
X8  & Negation scope is correct & \\
X9  & Boolean precedence from the request is preserved & \\
X10 & Grouping dimensions match the request & \\
X11 & Ordering does not distort interpretation of the result & \\
X12 & The denominator definition matches the indicator & \checkmark\\
X13 & The numerator definition matches the indicator & \checkmark\\
X14 & Compound institutional definitions are fully expanded & \checkmark\\
X15 & Deduplication level matches the counting object & \\
X16 & An ambiguous criterion is surfaced as TO-CONFIRM, not resolved silently & \checkmark\\
X17 & The TO-CONFIRM item states the competing candidates and what separates them & \checkmark\\
X18 & No condition appears in the query that is absent from the request & \checkmark\\
X19 & No condition appears in the request that is absent from the query & \checkmark\\
X20 & The indicator version in force on the reporting date is used & \\
X21 & Cross-campus scope matches the request & \\
X22 & A result-plausibility check is performed and reported & \\
\end{longtable}
\endgroup

\vspace{-6pt}
\noindent{\footnotesize Subset column: \checkmark\ marks the 18 items forming the
restricted instrument for prevalidated-criterion queries (C1--C3, C22--C24 =
field identifier correctness; X1--X3, X6, X7, X12--X14, X16--X19 = criterion
inclusion accuracy).}

%==============================================================================
\section{Scoring worksheet}
%==============================================================================
\noindent
For each query $q$ and expert $e$, with $A_{q}$ the set of applicable items in
dimension $D$:
\[
\text{score}_{D}(q,e)=\frac{1}{|A_{q}\cap D|}\sum_{i\in A_{q}\cap D} r_{i}(q,e),
\qquad
\text{score}_{D}(q)=\tfrac{1}{2}\bigl(\text{score}_{D}(q,1)+\text{score}_{D}(q,2)\bigr).
\]
Reported dimension values are the mean $\pm$ SD of $\text{score}_{D}(q)$ over the
evaluated queries. The LLM self-evaluation column was produced by issuing the
same 72 items to the generating model as a checklist, with no access to the
expert ratings.

\vspace{6pt}
\noindent
\PH{Attach the completed per-query rating matrix (queries $\times$ 72 items
$\times$ 2 raters) exported from the annotation tool on the compute cluster, and
report the observed Cohen $\kappa$ per dimension alongside the pooled value.}

\end{document}
